% Options for packages loaded elsewhere
\PassOptionsToPackage{unicode}{hyperref}
\PassOptionsToPackage{hyphens}{url}
%
\documentclass[
]{book}
\usepackage{amsmath,amssymb}
\usepackage{lmodern}
\usepackage{ifxetex,ifluatex}
\ifnum 0\ifxetex 1\fi\ifluatex 1\fi=0 % if pdftex
  \usepackage[T1]{fontenc}
  \usepackage[utf8]{inputenc}
  \usepackage{textcomp} % provide euro and other symbols
\else % if luatex or xetex
  \usepackage{unicode-math}
  \defaultfontfeatures{Scale=MatchLowercase}
  \defaultfontfeatures[\rmfamily]{Ligatures=TeX,Scale=1}
\fi
% Use upquote if available, for straight quotes in verbatim environments
\IfFileExists{upquote.sty}{\usepackage{upquote}}{}
\IfFileExists{microtype.sty}{% use microtype if available
  \usepackage[]{microtype}
  \UseMicrotypeSet[protrusion]{basicmath} % disable protrusion for tt fonts
}{}
\makeatletter
\@ifundefined{KOMAClassName}{% if non-KOMA class
  \IfFileExists{parskip.sty}{%
    \usepackage{parskip}
  }{% else
    \setlength{\parindent}{0pt}
    \setlength{\parskip}{6pt plus 2pt minus 1pt}}
}{% if KOMA class
  \KOMAoptions{parskip=half}}
\makeatother
\usepackage{xcolor}
\IfFileExists{xurl.sty}{\usepackage{xurl}}{} % add URL line breaks if available
\IfFileExists{bookmark.sty}{\usepackage{bookmark}}{\usepackage{hyperref}}
\hypersetup{
  pdftitle={Modelado y Análisis Espacial},
  pdfauthor={Gerardo Martín},
  hidelinks,
  pdfcreator={LaTeX via pandoc}}
\urlstyle{same} % disable monospaced font for URLs
\usepackage{longtable,booktabs,array}
\usepackage{calc} % for calculating minipage widths
% Correct order of tables after \paragraph or \subparagraph
\usepackage{etoolbox}
\makeatletter
\patchcmd\longtable{\par}{\if@noskipsec\mbox{}\fi\par}{}{}
\makeatother
% Allow footnotes in longtable head/foot
\IfFileExists{footnotehyper.sty}{\usepackage{footnotehyper}}{\usepackage{footnote}}
\makesavenoteenv{longtable}
\usepackage{graphicx}
\makeatletter
\def\maxwidth{\ifdim\Gin@nat@width>\linewidth\linewidth\else\Gin@nat@width\fi}
\def\maxheight{\ifdim\Gin@nat@height>\textheight\textheight\else\Gin@nat@height\fi}
\makeatother
% Scale images if necessary, so that they will not overflow the page
% margins by default, and it is still possible to overwrite the defaults
% using explicit options in \includegraphics[width, height, ...]{}
\setkeys{Gin}{width=\maxwidth,height=\maxheight,keepaspectratio}
% Set default figure placement to htbp
\makeatletter
\def\fps@figure{htbp}
\makeatother
\setlength{\emergencystretch}{3em} % prevent overfull lines
\providecommand{\tightlist}{%
  \setlength{\itemsep}{0pt}\setlength{\parskip}{0pt}}
\setcounter{secnumdepth}{5}
\usepackage{booktabs}
\ifluatex
  \usepackage{selnolig}  % disable illegal ligatures
\fi
\usepackage[]{natbib}
\bibliographystyle{apalike}

\title{Modelado y Análisis Espacial}
\author{Gerardo Martín}
\date{2022-01-11}

\begin{document}
\maketitle

{
\setcounter{tocdepth}{1}
\tableofcontents
}
\hypertarget{preuxe1mbulo}{%
\chapter{Preámbulo}\label{preuxe1mbulo}}

\begin{center}\includegraphics[width=6.03in]{logo} \end{center}

En el curso \textbf{Modelado y Análisis Espacial} aprenderemos a utilizar algunas herramientas para aprender a usar herramientas geográficas para analizar y representar procesos ambientales en el espacio. Los contenidos del índice se apegan al \href{Programa-curso.pdf}{programa completo del curso}, el cual se impartirá en los \href{Horario.pdf}{horarios normales establecidos}. Para conocer cuándo, cómo y qué temas se se impartirán puedes consultar la \href{Estrategia-docente.pdf}{estrategia docente}.

\hypertarget{encuadre-de-la-materia}{%
\chapter{Encuadre de la materia}\label{encuadre-de-la-materia}}

\hypertarget{criterios-de-evaluaciuxf3n}{%
\section{Criterios de evaluación}\label{criterios-de-evaluaciuxf3n}}

Las constribuciones a cada calificación parcial serán:

\begin{itemize}
\tightlist
\item
  Asistencia (25\%)
\item
  Trabajos de clase cumplidos (50\%)
\item
  Examen (25\%)
\item
  Participación (2 puntos extra máximo)
\end{itemize}

Cabe señalar, que la asistencia no corresponderá con su presencia en las sesiones sincrónicas, sino con el cumplimiento de los trabajos de clase. La participación se medirá tanto por participación directa en las sesiones sincrónicas como por el seguimiento que uds den a la clase por correo electrónico.

\hypertarget{cuxf3mo-se-daruxe1n-las-clases}{%
\section{¿Cómo se darán las clases?}\label{cuxf3mo-se-daruxe1n-las-clases}}

Trataré de apegarnos a los tiempos de actividades sincrónicas y asincrónicas establecidos en la \href{Estrategia-docente.pdf}{estrategia docente}, pero éstos son completamente flexibles. Una estrategia que me ha funcionado en cursos anteriores ha sido la de dedicar la última actividad sicrónica de la semana a resolver dudas, es en estas sesiones en las que hay muchas posibilidades de conseguir puntos de participación.

Todos los contenidos del curso, lecturas y presentaciones, se irán añadiendo a este sitio web conforme avanza el semestre. En el Google Classroom de la materia se irán anunciando las diferentes actividades y sesiones sincrónicas con anticipación suficiente. Igualmente, los examenes y resultados serán publicados a través de esta plataforma. A petición de uds, también se publicarán aquí los videos de las sesiones sincrónicas que tengamos, en especial para aquellos temas que sean mayor interés/dificultad/importancia. Los trabajos de práctica también se publicarán en Classroom.

Para finalizar, las clases sincrónicas estarán almacenadas en la playlist de la clase contenida en mi canal de youtube.

\hypertarget{reglas-del-saluxf3n}{%
\section{Reglas del salón}\label{reglas-del-saluxf3n}}

Estas, obviamente, son particulares del modelo en línea, por lo tanto aquí van las reglas de zoom:

\begin{enumerate}
\def\labelenumi{\arabic{enumi}.}
\tightlist
\item
  Micrófonos apagados
\item
  Cámaras prendidas, exceptuando:
  2.1. Si su velocidad de internet lo dificulta
  2.2. Si tienen datos limitados
\item
  Hacer muchas preguntas
\item
  Decirme si paso algo por alto
\end{enumerate}

\hypertarget{contacto}{%
\section{Contacto}\label{contacto}}

Para reportar fallos, resolver dudas y peticiones especiales grupales o individuales por favor enviar correo electrónico a \href{mailto:gerardo.mmc@enesmerida.unam.mx}{\nolinkurl{gerardo.mmc@enesmerida.unam.mx}}.

\hypertarget{unidad-i-introducciuxf3n-al-modelado-espacial}{%
\chapter{Unidad I: Introducción al modelado espacial}\label{unidad-i-introducciuxf3n-al-modelado-espacial}}

\hypertarget{unidad-ii}{%
\chapter{Unidad II}\label{unidad-ii}}

\hypertarget{unidad-iii}{%
\chapter{Unidad III}\label{unidad-iii}}

\hypertarget{unidad-iv}{%
\chapter{Unidad IV}\label{unidad-iv}}

\hypertarget{unidad-v}{%
\chapter{Unidad V}\label{unidad-v}}

\hypertarget{unidad-vi}{%
\chapter{Unidad VI}\label{unidad-vi}}

  \bibliography{book.bib,packages.bib,citas.bib}

\end{document}
